\chapter{Blockchains for Traders}\label{ch:m3:blockchains-for-traders}

A trader sees a \omterm{def:m3:blockchains-for-traders:stable}{stablecoin} priced at a discount on a decentralised venue and a premium on an
exchange. She withdraws \omterm{def:m3:blockchains-for-traders:key}{tokens} from the exchange to capture the gap. The withdrawal waits in a
queue, is broadcast, waits for a block, then for the confirmations the receiving venue requires;
the transaction's fee, bid in an auction she does not see, doubles while it waits. When the \omterm{def:m3:blockchains-for-traders:key}{tokens}
arrive the gap has closed. Crypto markets settle on \omterm{def:m3:blockchains-for-traders:chain}{blockchains}, and every property of a chain
that an engineer might find abstract (how blocks are made, when they are final, what a transaction
costs) is, for a trader, a latency, a cost or a risk. This chapter explains what a trader needs of
the machinery: ledgers and consensus, \omterm{def:m3:blockchains-for-traders:finality}{finality}, transactions and their fees, the \omterm{def:m3:blockchains-for-traders:mempool}{mempool}, the
layered chains built on top, \omterm{def:m3:blockchains-for-traders:bridge}{bridges} between them, and \omterm{def:m3:blockchains-for-traders:stable}{stablecoins}, the dollars of this market.

\section{Ledgers, consensus and finality}

\begin{definition}[Blockchain, consensus protocol]\label{def:m3:blockchains-for-traders:chain}
A \emph{blockchain}\index{blockchain} is a ledger kept by many independent computers as a chain of
blocks, each containing a batch of transactions and a cryptographic hash of the block before, so that
changing any past block would change every later one. A \emph{consensus protocol}\index{consensus
protocol} is the rule by which those computers agree on which block extends the chain.
\end{definition}

\begin{definition}[Proof of work, proof of stake, validator]\label{def:m3:blockchains-for-traders:pow}
Under \emph{proof of work}\index{proof of work} the right to propose a block is won by solving a
costly computational puzzle, and the chain with the most accumulated work is the valid one. Under
\emph{proof of stake}\index{proof of stake} it is assigned among \emph{validators}\index{validator},
participants who lock (stake) the chain's native \omterm{def:m3:blockchains-for-traders:key}{token} as collateral that can be destroyed if they
misbehave, and who propose and attest to blocks.
\end{definition}

\begin{definition}[Finality]\label{def:m3:blockchains-for-traders:finality}
\emph{Finality}\index{finality} is the point after which a transaction can no longer be reversed:
probabilistic under \omterm{def:m3:blockchains-for-traders:pow}{proof of work}, where each further block makes a reversal exponentially less likely;
economic under proof-of-stake protocols that finalise checkpoints, where reversing one would destroy a
large share of the stake.
\end{definition}

For a trader, \omterm{def:m3:blockchains-for-traders:finality}{finality} is a settlement cycle (One Quant Book 1, chapter 5) with a different shape. A
venue credits a deposit only after a number of confirmations it chooses, trading off its own risk of a
reversed deposit against its customers' wait.

\begin{dated}{2026-09}{Block times and finality}\label{dat:m3:blockchains-for-traders:times}
Bitcoin targets one block every ten minutes; a common convention treats six confirmations, about an hour,
as final. Ethereum moved from \omterm{def:m3:blockchains-for-traders:pow}{proof of work} to \omterm{def:m3:blockchains-for-traders:pow}{proof of stake} on 15 September 2022 (the Merge); time is
divided into 12-second slots and 32-slot epochs, and a block is finalised after two epochs, about 12.8
minutes. Ethereum block 26{,}048{,}449, produced at 16:17 UTC on 24 September 2026, had a \omterm{def:m3:blockchains-for-traders:gas}{gas} limit of 60 million,
of which it used 22.1 million, at a \omterm{def:m3:blockchains-for-traders:gas}{base fee} of 1.49 gwei.
\end{dated}

\section{Transactions, gas and the fee market}

\begin{definition}[Private key, token]\label{def:m3:blockchains-for-traders:key}
A \emph{private key}\index{private key} is the secret number whose signature authorises transactions from
an address; whoever holds it controls the address's assets. A \emph{token}\index{token} is a transferable
unit recorded by a \omterm{def:m3:blockchains-for-traders:chain}{blockchain}: the chain's native coin, or a unit whose balances a program on the chain
keeps.
\end{definition}

\begin{definition}[Smart contract, oracle]\label{def:m3:blockchains-for-traders:contract}
A \emph{smart contract}\index{smart contract} is a program deployed on a \omterm{def:m3:blockchains-for-traders:chain}{blockchain} whose code and state
are public and whose functions anyone can call by sending a transaction; the chain executes it and records
the result. An \emph{oracle}\index{oracle} is a mechanism that brings data from outside a chain (a price,
an event's outcome) into a smart contract.
\end{definition}

\begin{definition}[Gas, base fee, priority fee]\label{def:m3:blockchains-for-traders:gas}
\emph{Gas}\index{gas} is the unit in which Ethereum measures the computation and storage a transaction
uses; each block has a gas limit. A transaction pays, per unit of gas, the block's \emph{base
fee}\index{base fee}, which is burned and set by protocol, plus a \emph{priority fee}\index{priority fee}
(a tip) that it chooses and that goes to the block's proposer.
\end{definition}

\begin{proposition}[The base fee rule]\label{prop:m3:blockchains-for-traders:basefee}
Under EIP-1559 a block's \omterm{def:m3:blockchains-for-traders:gas}{gas} target is half its limit, and the next block's \omterm{def:m3:blockchains-for-traders:gas}{base fee} is the parent's
multiplied by $1 + \tfrac18\,(g - g^{*})/g^{*}$, where $g$ is the parent's \omterm{def:m3:blockchains-for-traders:gas}{gas} used and $g^{*}$ its target:
at most $+12.5\%$ after a full block and $-12.5\%$ after an empty one. After $k$ full blocks the \omterm{def:m3:blockchains-for-traders:gas}{base fee} is
$1.125^k$ times its starting level.
\end{proposition}

\begin{proof}
This is the EIP's update with $g = 2g^{*}$ (full) or $g = 0$ (empty); iterate for $k$ blocks. The integer
arithmetic of the protocol rounds each step down (and makes an increase at least one wei).
\end{proof}

Twenty full blocks, four minutes on Ethereum, multiply the \omterm{def:m3:blockchains-for-traders:gas}{base fee} by 10.5: from 10 gwei to 105.45 gwei. A
swap using 150{,}000 \omterm{def:m3:blockchains-for-traders:gas}{gas} with a 2-gwei tip then costs \$48.35 at an ether price of \$3{,}000 (illustrative),
against \$5.40 before. Fees are what a trader's latency is measured in when blocks are congested.

\begin{omfigure}
\begin{tikzpicture}
\begin{axis}[width=11cm,height=4.8cm,axis y line*=left,
  xlabel={block},ylabel={base fee, gwei},
  tick label style={font=\small},label style={font=\small},
  xmin=0,xmax=59,ymin=0,ymax=45,grid=major,grid style={gray!25},
  table/col sep=comma]
\addplot[omThm,very thick] table[x=block,y=base]{figdata/markets-3/14-blockchains-for-traders/shock.csv};\label{plot:m3:basefee}
\end{axis}
\begin{axis}[width=11cm,height=4.8cm,axis y line*=right,axis x line=none,
  ylabel={gas used, \% of limit},ylabel near ticks,
  tick label style={font=\small},label style={font=\small},
  xmin=0,xmax=59,ymin=0,ymax=110,
  legend columns=2,legend style={font=\footnotesize,at={(0.5,-0.3)},anchor=north,draw=none,fill=none,/tikz/every even column/.append style={column sep=8pt}},
  table/col sep=comma]
\addlegendimage{/pgfplots/refstyle=plot:m3:basefee}\addlegendentry{base fee (left)}
\addplot[omDef,thick,dashed,const plot] table[x=block,y=share]{figdata/markets-3/14-blockchains-for-traders/shock.csv};\addlegendentry{gas used (right)}
\end{axis}
\end{tikzpicture}

\omcaption{A demand shock under EIP-1559: from block 5 to block 34 demand for \omterm{def:m3:blockchains-for-traders:gas}{gas} is four times higher. Blocks
fill, the \omterm{def:m3:blockchains-for-traders:gas}{base fee} rises 12.5\% a block until demand falls back to the target (half the limit), and decays
when the shock ends. Illustrative demand; exact protocol arithmetic. Data: the chapter's tutorial.}\label{fig:m3:blockchains-for-traders:shock}
\end{omfigure}

\section{The mempool}

\begin{definition}[Mempool]\label{def:m3:blockchains-for-traders:mempool}
The \emph{mempool}\index{mempool} is the set of signed transactions that nodes have received and relayed
but that are not yet in a block: public, in most chains, to anyone who runs a node.
\end{definition}

A transaction in the public \omterm{def:m3:blockchains-for-traders:mempool}{mempool} is an order that everyone can see before it executes. That visibility
is the root of the extraction of \cref{ch:m3:maximal-extractable-value}: a searcher who sees a large swap
waiting can place its own transactions before and after it. It also means a trader's transfer is visible to
competitors, and that its inclusion depends on what others bid: the \omterm{def:m3:blockchains-for-traders:gas}{priority fee} is an auction for position,
and a transaction with too low a tip can wait for many blocks.

\begin{omfigure}
\centering
\begin{tikzpicture}[font=\small,
  st/.style={draw,thick,rounded corners=2pt,align=center,minimum height=1cm,text width=2.6cm},
  arr/.style={-{Stealth[length=2.2mm]},thick}]
\node[st,draw=omDef,fill=omDef!8] (w) at (0,0) {wallet signs\\with private key};
\node[st,draw=omProp,fill=omProp!8] (m) at (3.4,0) {public mempool\\(seen by all)};
\node[st,draw=omThm,fill=omThm!8] (b) at (6.8,0) {included in a\\block};
\node[st,draw=omExo,fill=omExo!8] (f) at (10.2,0) {confirmations,\\then finality};
\draw[arr] (w) -- (m); \draw[arr] (m) -- (b); \draw[arr] (b) -- (f);
\node[font=\footnotesize,anchor=north,align=center] at (5.1,-0.8) {seconds to minutes:\\tips and congestion};
\node[font=\footnotesize,anchor=north,align=center] at (10.2,-0.8) {minutes to an hour:\\the protocol};
\end{tikzpicture}

\omcaption{The life of a transaction and where a trader's time goes: waiting for inclusion depends on the
fee paid against others', waiting for \omterm{def:m3:blockchains-for-traders:finality}{finality} on the chain. Schematic.}\label{fig:m3:blockchains-for-traders:life}
\end{omfigure}

\section{Base layers, rollups, sequencers and bridges}

Base layers are slow and expensive because every node executes every transaction. Most activity has moved
to chains built on top of them.

\begin{definition}[Rollup, sequencer]\label{def:m3:blockchains-for-traders:rollup}
A \emph{rollup}\index{rollup} is a chain that executes transactions off the base layer and posts them, in
compressed batches, to the base layer, which then guarantees their data and, through fraud or validity
proofs, their correctness. Its \emph{sequencer}\index{sequencer} is the operator that orders the rollup's
transactions and produces its blocks; on most rollups it is a single operator.
\end{definition}

\begin{omfigure}
\centering
\begin{tikzpicture}[font=\small,
  l/.style={draw,thick,rounded corners=2pt,align=center,minimum height=0.9cm,text width=3.2cm},
  arr/.style={-{Stealth[length=2.2mm]},thick}]
\node[l,draw=omThm,fill=omThm!8] (seq) at (0,1.6) {rollup sequencer:\\orders transactions};
\node[l,draw=omDef,fill=omDef!8] (l1) at (0,-0.4) {base layer: stores batches,\\checks proofs};
\node[l,draw=omProp,fill=omProp!8] (other) at (7.2,1.6) {another chain};
\node[l,draw=omExo,fill=omExo!8] (br) at (7.2,-0.4) {bridge contracts:\\lock here, mint there};
\draw[arr] (seq) -- node[right,font=\footnotesize] {compressed batches} (l1);
\draw[arr] (l1) -- node[above,font=\footnotesize] {lock tokens} (br);
\draw[arr] (br) -- node[right,font=\footnotesize] {mint claim} (other);
\node[font=\footnotesize,anchor=east,align=right] at (-1.8,1.6) {users send\\transactions};
\end{tikzpicture}

\omcaption{A \omterm{def:m3:blockchains-for-traders:rollup}{rollup} orders transactions off the base layer and posts them there in batches; a \omterm{def:m3:blockchains-for-traders:bridge}{bridge} locks \omterm{def:m3:blockchains-for-traders:key}{tokens}
on one chain and mints a claim on them on another. The \omterm{def:m3:blockchains-for-traders:rollup}{sequencer} and the \omterm{def:m3:blockchains-for-traders:bridge}{bridge}'s contracts or signers are the
points of trust. Schematic.}\label{fig:m3:blockchains-for-traders:rollup}
\end{omfigure}

A single \omterm{def:m3:blockchains-for-traders:rollup}{sequencer} decides the order of transactions on its \omterm{def:m3:blockchains-for-traders:rollup}{rollup}, which is the same power, and the same
temptation, as a block builder's on the base layer; it also sees transactions before they are ordered. \omterm{def:m3:blockchains-for-traders:rollup}{Rollup}
transactions are cheap and fast to include, but they are final on the base layer only when the batch is.

\begin{dated}{2026-09}{Who sequences the rollups}\label{dat:m3:blockchains-for-traders:l2beat}
Of the 101 layer-2 projects in L2BEAT's summary on 24 September 2026, 3 had a decentralised set of \omterm{def:m3:blockchains-for-traders:rollup}{sequencers}; for 67, a user
could sequence his own transactions through the base layer if the \omterm{def:m3:blockchains-for-traders:rollup}{sequencer} failed, and for 21 there was no such mechanism.
\end{dated}

\begin{definition}[Bridge]\label{def:m3:blockchains-for-traders:bridge}
A \emph{bridge}\index{bridge} moves \omterm{def:m3:blockchains-for-traders:key}{tokens} between chains: it locks or burns them on one chain and releases
or mints a corresponding \omterm{def:m3:blockchains-for-traders:key}{token} on the other, on the strength of a proof or of a set of signers that attest
to the lock.
\end{definition}

A \omterm{def:m3:blockchains-for-traders:bridge}{bridge} concentrates assets in one set of contracts and keys, which makes it the largest target in crypto.
In March 2022 attackers who had obtained a majority of the signing keys of the Ronin \omterm{def:m3:blockchains-for-traders:bridge}{bridge} withdrew about
\$600 million of \omterm{def:m3:blockchains-for-traders:key}{tokens}; in April the US Treasury attributed the theft to a North Korean group.

\section{Stablecoins: mint and redeem}

\begin{definition}[Stablecoin, depeg]\label{def:m3:blockchains-for-traders:stable}
A \emph{stablecoin}\index{stablecoin} is a \omterm{def:m3:blockchains-for-traders:key}{token} designed to hold a fixed value against a currency, usually
the dollar, most often by an issuer that holds reserves (cash, bills, repo) and mints \omterm{def:m3:blockchains-for-traders:key}{tokens} against deposits
and burns them against redemptions. A \emph{depeg}\index{depeg} is a period in which the \omterm{def:m3:blockchains-for-traders:key}{token} trades away
from its target value in the secondary market.
\end{definition}

The peg is held by arbitrage between two markets. In the primary market, a few institutions deposit dollars
with the issuer and receive \omterm{def:m3:blockchains-for-traders:key}{tokens}, or return \omterm{def:m3:blockchains-for-traders:key}{tokens} and receive dollars, at one for one. In the secondary
market, everyone else trades \omterm{def:m3:blockchains-for-traders:key}{tokens} on exchanges. When the \omterm{def:m3:blockchains-for-traders:key}{token} trades below a dollar, an institution with
primary access buys it and redeems it; when it trades above, it mints and sells. The peg is only as good as
that arbitrage, which depends on the reserves being worth a dollar and on redemption being open.

\begin{omfigure}
\centering
\begin{tikzpicture}[font=\small,
  p/.style={draw,thick,rounded corners=2pt,align=center,minimum height=1cm,text width=2.5cm},
  arr/.style={-{Stealth[length=2.2mm]},thick}]
\node[p,draw=omDef,fill=omDef!8] (iss) at (0,0) {issuer\\(reserves)};
\node[p,draw=omThm,fill=omThm!8] (inst) at (5.2,0) {institution with\\primary access};
\node[p,draw=omProp,fill=omProp!8] (sec) at (10.4,0) {secondary market\\(exchanges, pools)};
\draw[arr] ([yshift=5pt]inst.west) -- node[above,font=\footnotesize] {\$ in, mint} ([yshift=5pt]iss.east);
\draw[arr] ([yshift=-5pt]iss.east) -- node[below,font=\footnotesize] {burn, \$ out} ([yshift=-5pt]inst.west);
\draw[arr,{Stealth[length=2.2mm]}-{Stealth[length=2.2mm]}] (inst) -- node[above,font=\footnotesize] {buys $<\$1$} node[below,font=\footnotesize] {sells $>\$1$} (sec);
\end{tikzpicture}

\omcaption{How a reserve-backed \omterm{def:m3:blockchains-for-traders:stable}{stablecoin} holds its peg: institutions with primary access mint and redeem at
one for one and trade the difference against the secondary market. If redemption stops or the reserves are
doubted, nothing holds the peg. Schematic.}\label{fig:m3:blockchains-for-traders:mint}
\end{omfigure}

\begin{dated}{2026-09}{Stablecoins}\label{dat:m3:blockchains-for-traders:stable}
Total \omterm{def:m3:blockchains-for-traders:stable}{stablecoin} supply was about USD~303 billion on 10 September 2026, of which USDT about USD~183 billion and
USDC about USD~74 billion (aggregator figures). In March 2023 Circle could not withdraw USD~3.3 billion of USDC's
reserves, about 8\% of them, from the failed Silicon Valley Bank; USDC traded as low as about 86 cents before the
reserves were recovered and the peg restored. The Federal Reserve's researchers note that most holders cannot
redeem from the issuer and can only trade in secondary markets, and that issuance and redemption were constrained
by banking hours.
\end{dated}

\section{Tutorial: the base fee and the cost of a transaction}\label{tut:m3:blockchains-for-traders:gas}

\begin{tutorial}
\textbf{Goal.} Implement the base-fee rule in the protocol's integer arithmetic, project it through full
blocks and a demand shock, and cost a transaction. \textbf{End state:}
\cref{fig:m3:blockchains-for-traders:shock} and the numbers of the weekend problem.
\begin{tutsteps}
\item \textbf{The rule.} \Cref{prop:m3:blockchains-for-traders:basefee} as the EIP writes it.
\omcode{code/firm/gasfee/firm_gasfee.py}{12}{35}{The EIP-1559 base fee, its projection and a transaction's cost.}\label{lst:m3:blockchains-for-traders:basefee}
\item \textbf{The shock.} Demand for \omterm{def:m3:blockchains-for-traders:gas}{gas} falls with the \omterm{def:m3:blockchains-for-traders:gas}{base fee}; the fee rises until demand meets the target.
\omcode{code/markets-3/14-blockchains-for-traders/python/m3_chain.py}{20}{33}{A demand shock under the base-fee rule.}\label{lst:m3:blockchains-for-traders:shock}
\item \textbf{Run} \texttt{m3\_chain.full\_blocks(20)}, \texttt{arbitrage\_breakeven()} and \texttt{fig\_chain.py}.
\end{tutsteps}
\textbf{What to change next.} Make the shock's demand more elastic and see how high the \omterm{def:m3:blockchains-for-traders:gas}{base fee} goes; add a
distribution of competitors' tips and compute the tip that maximises an arbitrage's expected profit.
\end{tutorial}

\section{Build: fees and confirmation times}\label{bld:m3:blockchains-for-traders:gasfee}

\begin{build}
\bfield{Purpose} Every on-chain action of the miniature firm (a transfer between venues, a swap, a liquidation)
must be costed and timed before it is sent: the fee component projects \omterm{def:m3:blockchains-for-traders:gas}{base fees} and prices transactions; the
confirmation model tells the transfer planner of \cref{ch:m3:spot-markets} how long funds are in flight.
\bfield{Interface} \texttt{next\_base\_fee(parent\_base\_fee, gas\_used, gas\_limit)} in wei;
\texttt{project(base\_fee, usage, gas\_limit)}; \texttt{tx\_cost\_usd(gas, base\_fee\_gwei, priority\_gwei,
eth\_usd)}; \texttt{confirmation\_seconds(chain, confirmations)}.
\bfield{Rules} Integer wei arithmetic exactly as in the EIP (floor division, minimum increase of one wei);
unknown chains are an error, never a default.
\bfield{Acceptance tests} \texttt{code/firm/gasfee/tests/}: $\pm12.5\%$ at full and empty blocks; unchanged at
target; the one-wei minimum; twenty full blocks give $1.125^{20}$; confirmation times.
\bfield{Stretch} Blob fees for \omterm{def:m3:blockchains-for-traders:rollup}{rollup} data; other chains' fee markets; inclusion probability as a function of the
tip from observed \omterm{def:m3:blockchains-for-traders:mempool}{mempools}.
\end{build}

\begin{omsources}
\item S.~Nakamoto, \emph{Bitcoin: A Peer-to-Peer Electronic Cash System}, 2008.
\item V.~Buterin et al., \emph{EIP-1559: Fee market change for ETH 1.0 chain}; ethereum.org, proof-of-stake and The
Merge documentation.
\item Ethereum block 26{,}048{,}449 via a public JSON-RPC node; L2BEAT, scaling summary, 24 September 2026.
\item Federal Reserve, FEDS Notes, ``In the shadow of bank runs: lessons from the Silicon Valley Bank failure and its
impact on stablecoins'', 17 December 2025.
\item CoinDesk, ``US officials tie North Korea's Lazarus hackers to \$625M Axie Infinity Ronin exploit'', 14 April
2022.
\item Stablecoin supply: StablecoinBeat tracker, September 2026.
\end{omsources}

\section{Exercises}

\begin{exercise}[$\star$]\label{exo:m3:blockchains-for-traders:1}
The \omterm{def:m3:blockchains-for-traders:gas}{base fee} is 20 gwei and the parent block used 75\% of its limit. Give the next \omterm{def:m3:blockchains-for-traders:gas}{base fee}.
\end{exercise}

\begin{exercise}[$\star$]\label{exo:m3:blockchains-for-traders:2}
How long does a venue that requires six Bitcoin confirmations make a depositor wait, on average? And for an
Ethereum deposit credited at \omterm{def:m3:blockchains-for-traders:finality}{finality}?
\end{exercise}

\begin{exercise}[$\star$]\label{exo:m3:blockchains-for-traders:3}
Why does a \omterm{def:m3:blockchains-for-traders:stable}{stablecoin}'s peg depend on who may redeem it?
\end{exercise}

\begin{exercise}[$\star\star$]\label{exo:m3:blockchains-for-traders:4}
How many consecutive full blocks double the \omterm{def:m3:blockchains-for-traders:gas}{base fee}? How many empty blocks halve it?
\end{exercise}

\begin{exercise}[$\star\star$]\label{exo:m3:blockchains-for-traders:5}
A swap uses 150{,}000 \omterm{def:m3:blockchains-for-traders:gas}{gas}; the \omterm{def:m3:blockchains-for-traders:gas}{base fee} is 30 gwei, the tip 2 gwei and ether \$3{,}000. What does it cost?
\end{exercise}

\begin{exercise}[$\star\star$]\label{exo:m3:blockchains-for-traders:6}
Explain why a \omterm{def:m3:blockchains-for-traders:rollup}{rollup} transaction can be cheap and fast and yet not final.
\end{exercise}

\begin{exercise}[$\star\star\star$]\label{exo:m3:blockchains-for-traders:7}
\emph{Coding.} With \texttt{demand\_shock}, find the highest \omterm{def:m3:blockchains-for-traders:gas}{base fee} reached and the block at which blocks first
stop being full.
\end{exercise}

\begin{exercise}[$\star\star\star$]\label{exo:m3:blockchains-for-traders:8}
\emph{Find the flaw.} ``A \omterm{def:m3:blockchains-for-traders:stable}{stablecoin} backed one for one by Treasury bills cannot lose its peg.''
\end{exercise}

\section{Problem: The Gas Spike}

\begin{problem}[{Weekend problem --- an arbitrage racing a congested chain}]\label{pb:m3:blockchains-for-traders:1}
A searcher runs an arbitrage that earns \$60 before \omterm{def:m3:blockchains-for-traders:gas}{gas}, uses 150{,}000 \omterm{def:m3:blockchains-for-traders:gas}{gas} and tips 2 gwei. The \omterm{def:m3:blockchains-for-traders:gas}{base fee} is 10 gwei and
ether costs \$3{,}000 (illustrative). A popular \omterm{def:m3:blockchains-for-traders:key}{token} launch fills every block.

\medskip\noindent\textbf{Part I --- The fee.}
\begin{enumerate}
\item What does the arbitrage pay in \omterm{def:m3:blockchains-for-traders:gas}{gas} now?
\item What is the \omterm{def:m3:blockchains-for-traders:gas}{base fee} after one full block, and after twenty?
\item What does the arbitrage then pay in \omterm{def:m3:blockchains-for-traders:gas}{gas}?
\item How long do twenty blocks take?
\item Who receives the \omterm{def:m3:blockchains-for-traders:gas}{base fee}, and who the tip?
\end{enumerate}

\medskip\noindent\textbf{Part II --- The break-even.}
\begin{enumerate}[resume]
\item At what \omterm{def:m3:blockchains-for-traders:gas}{base fee} does the arbitrage stop paying?
\item After how many full blocks is that reached?
\item What if the arbitrage earned \$120?
\item Why might the searcher raise its tip during the spike?
\item What happens to the \omterm{def:m3:blockchains-for-traders:gas}{base fee} when the launch ends?
\end{enumerate}

\medskip\noindent\textbf{Part III --- Alternatives.}
\begin{enumerate}[resume]
\item Would a \omterm{def:m3:blockchains-for-traders:rollup}{rollup} be cheaper, and what would the searcher give up?
\item How does a private submission change the searcher's competition?
\item Why is the \omterm{def:m3:blockchains-for-traders:gas}{base fee} burned rather than paid to the proposer?
\item What does the spike do to users who are not arbitrageurs?
\item What does the searcher learn from the \omterm{def:m3:blockchains-for-traders:mempool}{mempool} during the spike?
\end{enumerate}

\medskip\noindent\textbf{Part IV --- Judgement.}
\begin{enumerate}[resume]
\item Is the fee market a good way to allocate block space?
\item What does a trader moving funds between venues need from this chapter?
\item Why do fees matter more for small arbitrages than large ones?
\item State the \emph{named result}: the base-fee multiplier after twenty full blocks, and the number of full blocks
after which the arbitrage stops paying.
\item In one sentence: what does a transaction's fee buy?
\end{enumerate}
\end{problem}

\section{Interview questions}

\begin{interviewq}[$\star$ \iqroles{trader}]\label{iq:m3:blockchains-for-traders:1}
What does \omterm{def:m3:blockchains-for-traders:finality}{finality} mean, and why does it matter for moving funds between exchanges?
\end{interviewq}

\begin{interviewq}[$\star$ \iqroles{developer}]\label{iq:m3:blockchains-for-traders:2}
Explain EIP-1559's \omterm{def:m3:blockchains-for-traders:gas}{base fee} and \omterm{def:m3:blockchains-for-traders:gas}{priority fee}.
\end{interviewq}

\begin{interviewq}[$\star\star$ \iqroles{researcher, trader}]\label{iq:m3:blockchains-for-traders:3}
How does a reserve-backed \omterm{def:m3:blockchains-for-traders:stable}{stablecoin} keep its peg, and how can it fail?
\end{interviewq}

\begin{interviewq}[$\star\star$ \iqroles{risk}]\label{iq:m3:blockchains-for-traders:4}
What are the risks of holding assets bridged from one chain to another?
\end{interviewq}

\begin{interviewq}[$\star\star$ \iqroles{developer}]\label{iq:m3:blockchains-for-traders:5}
A deposit shows one confirmation. When should your system credit it, and what can go wrong?
\end{interviewq}

\begin{interviewq}[$\star\star\star$ \iqroles{developer, researcher}]\label{iq:m3:blockchains-for-traders:6}
Design a fee estimator that tells a trading system what tip to pay to be included within two blocks with 95\%
probability.
\end{interviewq}
